\section{逆过程（Reverse Process）}

\subsection{逆过程的定义}

逆过程（Reverse Process）是生成过程的核心——从纯噪声 $x_T \sim \mathcal{N}(0, I)$ 出发，逐步去噪直到得到数据样本 $x_0$。

逆过程的联合分布定义为：

$$p_\theta(x_0, x_1, \ldots, x_T) = p(x_T) \prod_{t=1}^{T} p_\theta(x_{t-1} | x_t)$$

其中每一步的逆转移分布参数化为高斯：

$$p_\theta(x_{t-1} | x_t) = \mathcal{N}(x_{t-1}; \mu_\theta(x_t, t), \sigma_t^2 I)$$

\begin{importantbox}{逆过程的物理含义}
逆过程是前向过程的``时间反演''。前向过程逐步加噪，逆过程逐步去噪。神经网络 $\mu_\theta(x_t, t)$ 的任务是：给定当前的噪声图像 $x_t$ 和时间步 $t$，预测去噪一步后的均值。注意神经网络需要知道 $t$ 的信息，因为不同时间步的噪声水平不同。
\end{importantbox}

\subsection{逆过程的生成流程}

给定训练好的模型，生成过程如下：

\begin{enumerate}
    \item 从标准高斯分布采样 $x_T \sim \mathcal{N}(0, I)$
    \item 对于 $t = T, T-1, \ldots, 1$：
    \begin{itemize}
        \item 使用神经网络计算 $\mu_\theta(x_t, t)$ 和 $\sigma_t$
        \item 采样 $x_{t-1} \sim \mathcal{N}(\mu_\theta(x_t, t), \sigma_t^2 I)$
    \end{itemize}
    \item 最终得到生成样本 $x_0$
\end{enumerate}

\begin{warningbox}{生成速度问题}
逆过程需要执行 $T$ 步顺序采样（DDPM 中 $T = 1000$），每一步都需要一次神经网络前向传播。这使得 DDPM 的生成速度远慢于 GAN（仅需一次前向传播）。这也是后续 DDIM、一致性模型等工作要解决的核心问题。
\end{warningbox}

\subsection{本章小结}

逆过程是扩散模型中唯一需要学习的部分。通过参数化每一步的逆转移为高斯分布，并用神经网络预测其均值，模型学会了从噪声中逐步恢复数据。


